% STATUS: written 2026-11-05, recalled results as is (one environment each, source kind and whole statement)
% Sources read in full: Woronowicz 1989 (pp. 125--170, Section 1), Losert 2008 (pp. 221--246, Sections 1--2 and the statements of the main results).
\section{Crossed products and crossed homomorphisms}\label{sec:group-hopf:crossed-products}

Section~\ref{sec:group-hopf:quantum} ended with quotients of the kernel of the
multiplication map. The same tensor $1\otimes b-b\otimes 1$ has a group form, the
difference $\mu-x\circ\mu$ of Definition~\ref{def:group-hopf:crossed-hom} below.
When $G$ acts on the ground ring, the Leibniz rule restricted
to group elements reads $D(xy)=D(x)y+xD(y)$, and the map $x\mapsto D(x)x^{-1}$ of
Section~\ref{sec:group-hopf:group-algebras} becomes a crossed homomorphism. We
first recall the cocycle that builds a crossed product.

\begin{definition}[{\cite[second condition, p.~567]{Davydov2013}}]\label{def:group-hopf:cocycle}
An element $\varphi$ of $H\otimes H$ satisfies the \emph{2-cocycle condition} for
the comultiplication $\Delta$ of a bialgebra $H$ when
\[
1\otimes\varphi+(I\otimes\Delta)(\varphi)=\varphi\otimes1+(\Delta\otimes I)(\varphi).
\]
\end{definition}

The cocycle condition is the second equation of a twisted derivation,
Definition~\ref{def:group-hopf:twisted}; it says that the first components of
$\varphi$ in a coassociative expansion agree. Its group version is older: Losert
calls it a crossed homomorphism, with a left Banach $G$-module and measures in
place of tensors \cite[p.~222]{Losert2008}.

\begin{definition}[{\cite[p.~222]{Losert2008}}]\label{def:group-hopf:crossed-hom}
Let $\Omega$ be a locally compact space and $G$ a discrete group acting on $\Omega$
by homeomorphisms. A mapping $\Phi\colon G\to M(\Omega)$ (or more generally,
$\Phi\colon G\to X$, where $X$ is a left Banach $G$-module) is called a
\emph{crossed homomorphism} if $\Phi(xy)=\Phi(x)+x\circ\Phi(y)$ for all $x,y\in G$.
The map $\Phi$ is called \emph{bounded} if $\sup_{x\in G}\|\Phi(x)\|<\infty$. For
$\mu\in M(\Omega)$, the special example $\Phi_\mu(x)=\mu-x\circ\mu$ is called a
\emph{principal crossed homomorphism}. It is the principal derivation of
Section~\ref{sec:basic-theory:inner-outer} under its third name, and the three
are one object: an inner derivation of a ring is the same map as a principal
crossed homomorphism of a group.
\end{definition}

In $\Phi_\mu=\mu-x\circ\mu$ the measure $\mu$ plays the role of the inducing
element of an inner derivation, and Losert proves the converse for every group:
Theorem~\ref{thm:group-hopf:losert} below makes every bounded crossed homomorphism
principal. The derivation problem of Johnson asks whether every derivation
$D\colon L^1(G)\to L^1(G)$ is inner, induced by a bounded measure
\cite[p.~221]{Losert2008}. The question goes back to Williamson around 1965, and
Johnson settled the discrete groups, the amenable groups and the connected groups
\cite[p.~221]{Losert2008}. Losert proves the full statement through crossed
homomorphisms \cite[p.~222]{Losert2008}.

\begin{theorem}[{\cite[Theorem~1.1, p.~222]{Losert2008}}]\label{thm:group-hopf:losert}
Let $\Omega$ be a locally compact space, $G$ a discrete group with a left action of
$G$ on $\Omega$ by homeomorphisms. Then any bounded crossed homomorphism $\Phi$
from $G$ to $M(\Omega)$ is principal. There exists $\mu\in M(\Omega)$ with
$\|\mu\|\le 2\|\Phi\|$ such that $\Phi=\Phi_\mu$.
\end{theorem}

So a derivation $D\colon L^1(G)\to M(G)$, having values in $L^1(G)$ and being
bounded, defines a bounded crossed homomorphism from $G$ to $M(G)$ for the
conjugation action of $G$ on $G$; Theorem~\ref{thm:group-hopf:losert} makes it
principal, and the inducing measure gives the inner derivation
\cite[p.~222]{Losert2008}.

\begin{corollary}[{\cite[Corollary~1.2, p.~222]{Losert2008}}]\label{cor:group-hopf:losert}
Let $G$ denote a locally compact group. Then any derivation
$D\colon L^1(G)\to M(G)$ is inner.
\end{corollary}

So the chapter closes where it began: a derivation of a group algebra is again the
bracket with a single element, once the group and the norm allow it. From the
integral group ring of a finite group to the measure algebra of a locally compact
group, only three ways for a derivation to fail to be inner appeared. % claim-ok: synthesis of the cited results of this chapter; no new external fact
So they are the free factors of Theorem~\ref{thm:group-hopf:spiegel-abelian}, the
infinite sum of Example~\ref{ex:group-hopf:outer}, and the twists of
Section~\ref{sec:group-hopf:hopf}.
